\section{Conclusion}\label{sec:conclusion}

We developed a mathematical framework for misalignment sensitivity thresholds
in softmax-parameterized trading agents. The central object is the threshold
$\tau(\delta)$, the minimum perturbation intensity producing KL divergence
exceeding $\delta$. Our main theorem characterizes this threshold via the
Fisher information of the policy parameterization:
$\tau = \sqrt{2\delta/\lVert G \rVert_F^2} + O(\delta)$.

The principal finding is the temperature scaling law $\tau \propto T$
(Corollary~\ref{cor:temperature}), which identifies the softmax temperature
as the dominant determinant of misalignment robustness. This result, confirmed
computationally with $R^2 = 0.999$, has a clear practical interpretation:
agents with more concentrated decision distributions (lower $T$) are more
vulnerable to reward perturbations. Market volatility and decision horizon
are asymptotically irrelevant (Propositions~\ref{prop:volatility}
and~\ref{prop:horizon}), refuting the initially hypothesized power-law
dependence on these parameters.

These results redirect attention from external market conditions to internal
policy architecture. The threshold formula provides a concrete, computable
safety metric for deployment decisions: given a maximum plausible perturbation
intensity and a desired significance level, one can determine the minimum
temperature required for robust operation. Future work should extend the
framework to non-softmax parameterizations, multi-agent systems, and dynamic
perturbation models.
